\documentclass[11pt]{article}
\usepackage[a4paper,margin=27mm]{geometry}
\usepackage{amsmath,amssymb,amsthm,booktabs,array}
\usepackage[colorlinks=true,linkcolor=blue,urlcolor=blue,citecolor=blue]{hyperref}
\hypersetup{pdftitle={Degree forty-four in the closure of seven-product polynomial evaluations},pdfauthor={Marcus Webb}}
\newtheorem{theorem}{Theorem}
\newtheorem{lemma}[theorem]{Lemma}
\newtheorem{corollary}[theorem]{Corollary}
\newcommand{\C}{\mathbb C}
\newcommand{\V}[1]{\C[x]_{\le #1}}
\newcommand{\Span}{\operatorname{span}}
\newcommand{\cl}{\overline}
\title{Degree forty-four in the closure of seven-product polynomial evaluations\\
\large A negative resolution of MF-14}
\author{Marcus Webb\\The University of Manchester}
\date{17 September 2026}
\begin{document}
\maketitle
\begin{abstract}
Starting from $1,x$ and allowing free complex linear combinations, we prove
that the Zariski closure of the polynomials computable with seven products
contains every polynomial of degree at most $44$. This refutes the equality
to $42$ in MF-14. A simultaneous four-product degeneration retains two
intermediate polynomials as well as an arbitrary polynomial of degree at most
$12$. Three further products give a coefficient map with a $45\times45$
integer Jacobian of determinant $256$. The published dimension theorem of
Jarlebring and Lorentzon additionally implies $44\le d_7\le47$.
\end{abstract}

\noindent\textbf{Provenance and review.}
The construction and earlier computational checks were developed with
substantial ChatGPT assistance in the conversation supplied by the author
\cite{chat}. This standalone exposition uses the polynomial-determinant
version of the border argument in that conversation. Separate Codex agents
reviewed the border lemma and the continuation/certificate; their reports and
independently written exact verifiers accompany this manuscript. These are
informal AI-agent reviews, not external human peer review or formal verification.

\section{Statement and computational model}
Let $\mathcal P_m$ be the outputs obtainable from $1,x$ using arbitrary
complex linear combinations of previously available polynomials and at most
$m$ products of two previously available polynomials. All intermediate
polynomials have degree at most $2^m$. Put
\[
 X_7=\cl{\mathcal P_7}^{\,Z}\subseteq\V{128},\qquad
 d_7=\max\{d:\V d\subseteq X_7\}.
\]
All closures below are Zariski closures in the indicated full coefficient
spaces. In particular, approximating outputs are allowed to have degree
greater than the degree of their limits. No coefficient projection is an
allowed computational operation.

\begin{theorem}\label{thm:main}
One has $\V{44}\subseteq X_7$. Consequently the equality $d_7=42$ asked
in MF-14, and in Conjecture 13 of \cite{JL}, is false.
\end{theorem}

The previously audited degree-$42$ coverage result of Colbrook
\cite{Colbrook} remains valid. The new proof uses degenerating intermediate
computations. It does not assert exact representation of every degree-$44$
polynomial, real Euclidean density, or stable coefficient recovery.

\section{Simultaneous availability after four products}
Let $\mathcal A_4\subseteq(\V{16})^4$ be the ordered quadruples
simultaneously available after at most four products, including arbitrary
free linear combinations of the computed polynomials. Set
$\mathcal B_4=\cl{\mathcal A_4}^{\,Z}$.

\begin{lemma}\label{lem:border}
For all $\alpha,\beta\in\C$ and $P\in\V{12}$,
\[
 (x^2,x^4+\alpha x^3,x^5+\beta x^3,P)\in\mathcal B_4.
\]
\end{lemma}

\begin{proof}
Fix $\alpha,\beta$, put $\eta=\beta+\alpha^2$, and initially fix an
arbitrary $\gamma\in\C$. For nonzero $s$ tending to zero, put
\begin{equation}\label{eq:QR}
 Q=x^4+\alpha x^3,\qquad
 R_s=s^2Q^2+\gamma s x^2Q+xQ+\eta x^3.
\end{equation}

\paragraph{Three products.}
The first two products give $x^2$ and $Q=x^2(x^2+\alpha x)$.
For $a=s^2$, $b=\gamma s$, define
\[
 g=b-a,\quad \delta=b-2a,\quad
 \sigma=\frac{\eta-1+\alpha g}{\delta},\quad \rho=1-a\sigma.
\]
Direct multiplication gives the identity
\begin{equation}\label{eq:third}
 R_s=(aQ+\rho x+g x^2)(Q+\sigma x+x^2)-gQ-\rho\sigma x^2.
\end{equation}
Every factor uses only already available polynomials, so this costs one
additional product. The denominator $s(\gamma-2s)$ is nonzero for all
sufficiently small nonzero $s$, also when $\gamma=0$. Moreover,
\begin{equation}\label{eq:Rlimit}
 R_s-\alpha Q\longrightarrow R_0-\alpha Q=x^5+\beta x^3.
\end{equation}

\paragraph{The fourth-product image is dense in its product span.}
Let $V_s=\Span\{1,x,x^2,Q,R_s\}$ and
$W_s=\Span\{uv:u,v\in V_s\}$. Products from
$U=\Span\{1,x,x^2,Q\}$ span exactly
$\V6+\Span\{Q^2\}$: obtain $x^3,x^4$ directly and then
$x^5=xQ-\alpha x^4$, $x^6=x^2Q-\alpha x^5$.
Since $R_s$ already lies in that span,
\begin{equation}\label{eq:W}
 W_s=\V6+\Span\{Q^2,xR_s,x^2R_s,QR_s,R_s^2\}.
\end{equation}
The five extra polynomials have respective degrees $8,9,10,12,16$ and
nonzero leading coefficients for $s\ne0$. Thus $\dim W_s=12$.

Consider $\Psi_s:V_s^3\to W_s$, $(u,v,w)\mapsto uv+w$.
At $u=R_s$, $v=Q+\lambda x^2$, its differential includes these twelve columns:
\begin{equation}\label{eq:minorcols}
 1,x,x^2,Q,R_s,\ x(Q+\lambda x^2),x^2(Q+\lambda x^2),
 Q(Q+\lambda x^2),\ xR_s,x^2R_s,QR_s,R_s^2.
\end{equation}
In coefficient rows $0,1,2,3,4,5,6,8,9,10,12,16$, their determinant is
\begin{equation}\label{eq:minor}
 a^5\bigl(\lambda-\eta-\alpha\lambda(b-a\lambda)\bigr).
\end{equation}
This polynomial identity follows by coefficient expansion; the accompanying
exact symbolic verifier also recomputes it. Choose $\lambda=\eta+1$.
The determinant becomes
\[
 s^{10}\bigl[1-\alpha(\eta+1)(\gamma s-(\eta+1)s^2)\bigr],
\]
which is nonzero for sufficiently small nonzero $s$, for every fixed choice
of $\alpha,\eta,\gamma$. Hence $\Psi_s$ has full differential rank.
Restricting to twelve suitable affine directions and applying the complex
inverse function theorem shows that its image contains a Euclidean-open
subset of $W_s$, so it is Zariski dense in $W_s$.
All actual outputs $uv+w$ cost one further product and retain the previous
intermediates. Closedness therefore gives
\begin{equation}\label{eq:joint}
 (x^2,Q,R_s-\alpha Q,p)\in\mathcal B_4\qquad(p\in W_s).
\end{equation}
This is a closure argument, not an identification of a sum of products with
one actual product.

\paragraph{An explicit polynomial degeneration.}
Set $(f_1,f_2,f_3,f_4,f_5)=(x^2R_s,Q^2,QR_s,R_s^2,xR_s)$ and define
\begin{equation}\label{eq:E}
 E_s(x)=\det\begin{pmatrix}
 [x^7]f_1&\cdots&[x^7]f_5\\
 [x^8]f_1&\cdots&[x^8]f_5\\
 [x^9]f_1&\cdots&[x^9]f_5\\
 [x^{10}]f_1&\cdots&[x^{10}]f_5\\
 f_1(x)&\cdots&f_5(x)
 \end{pmatrix}.
\end{equation}
Expansion along the last row gives $E_s\in W_s$ and repeated rows give
$[x^j]E_s=0$ for $7\le j\le10$. Direct polynomial expansion gives
\begin{align}
 [x^{11}]E_s&=(4\gamma^2-3\alpha)s^4+O(s^5),\label{eq:lead11}\\
 [x^{12}]E_s&=-s^4+O(s^5),\label{eq:lead12}
\end{align}
while the coefficients in degrees $13,14,15,16$ are divisible by
$s^5,s^6,s^7,s^7$, respectively. Here the statements are identities in
$\mathbb Z[\alpha,\eta,\gamma,s]$, so the remainder notation denotes exact
divisibility. The verifier expands the cofactors of $E_s$ with integer
polynomial arithmetic and checks every coefficient and divisibility.
Appendix~\ref{app:cofactors} supplies explicit identities
for them. Degree is at most $16$.

Write $\pi_{\le6}$ for coefficient truncation used only in this algebraic
identity. Since $\V6\subseteq W_s$, the polynomial
\begin{equation}\label{eq:Z}
 Z_s=-s^{-4}(E_s-\pi_{\le6}E_s)\in W_s\quad(s\ne0)
\end{equation}
has coefficients polynomial in $s$, and satisfies
\begin{equation}\label{eq:Zlimit}
 Z_0=x^{12}+(3\alpha-4\gamma^2)x^{11}.
\end{equation}
The subtraction in \eqref{eq:Z} occurs inside $W_s$; no truncation of an
actual circuit output is being charged as free.

The eleven polynomials
\[
 1,x,\ldots,x^6,\quad x^2R_s,\quad Q^2,\quad QR_s,\quad R_s^2
\]
converge to a basis of $\V{10}$. Indeed $R_0=x^5+\alpha x^4+\eta x^3$,
and the four last limits are monic of degrees $7,8,9,10$; their coefficient
matrix together with the first seven monomials is triangular with determinant
one.

For $P=\sum_{j=0}^{12}p_jx^j$ with $p_{12}\ne0$, choose a complex root
\[
 \gamma^2=(3\alpha-p_{11}/p_{12})/4.
\]
Then $P-p_{12}Z_0\in\V{10}$. Express this difference in the limiting
eleven-element basis and use the same constant coefficients at $s\ne0$.
Adding $p_{12}Z_s$ gives $p_s\in W_s$ with $p_s\to P$ in all coefficients
through degree $16$. Taking limits in \eqref{eq:joint}, using
\eqref{eq:Rlimit} and Euclidean closedness of the algebraic set
$\mathcal B_4$, proves the lemma for $p_{12}\ne0$. For $p_{12}=0$, apply
that result to $P+\varepsilon x^{12}$ and take $\varepsilon\to0$.
\end{proof}

\section{Three further products and an exact coefficient certificate}
Take
\[
 Q=x^4+\alpha x^3,\quad R=x^5+\beta x^3,\quad
 P=x^{12}+\sum_{j\in J}\xi_jx^j,\qquad J=(3,6,7,8,9,10,11).
\]
Define $L_3=(x,x^2,Q)$, $L_4=(x,x^2,Q,R)$, $L_5=(x,x^2,Q,R,P)$,
and compute
\begin{align}
 F&=(P+\mathbf u\cdot L_4)(R+\mathbf v\cdot L_3),\label{eq:F}\\
 G&=(F+\mathbf a\cdot L_5)(R+\mathbf b\cdot L_3),\label{eq:G}\\
 H&=(G+\mathbf c\cdot L_6)(G+\mathbf d\cdot L_6),
 \qquad L_6=(x,x^2,Q,R,P,F).\label{eq:H}
\end{align}
The final output is
\begin{equation}\label{eq:output}
 p=\mathbf z\cdot(1,x,x^2,Q,R,P,F,G,H).
\end{equation}
The degrees of $F,G,H$ are at most $17,22,44$ for all parameters.
Order the parameters by
\begin{equation}\label{eq:parameters}
 (\alpha,\beta,\xi_3,\xi_6,\xi_7,\xi_8,\xi_9,\xi_{10},\xi_{11},
 \mathbf u,\mathbf v,\mathbf a,\mathbf b,\mathbf c,\mathbf d,\mathbf z),
\end{equation}
with entries within each vector in their displayed order. The vector lengths
are $4,3,5,3,6,6,9$, and $\mathbf z=(z_0,\ldots,z_8)$.
There are $2+7+4+3+5+3+6+6+9=45$ parameters, defining a polynomial map
$\Phi:\C^{45}\to\V{44}$ with integer coefficients.

\paragraph{Why its image lies in the full seven-product closure.}
Fix the continuation coefficients in \eqref{eq:F}--\eqref{eq:output}.
They define a polynomial map $T:(\V{16})^4\to\V{128}$ using the four
input entries in place of $x^2,Q,R,P$, while retaining the free input $x$.
On arbitrary inputs its product degrees are at most $32,48,96$, so its
codomain is indeed the full degree-$128$ space. For every quadruple in
$\mathcal A_4$, the output costs at most seven products. By continuity,
\[
 T(\mathcal B_4)\subseteq\cl{T(\mathcal A_4)}^{\,Z}\subseteq X_7.
\]
For example the first inclusion follows by taking the inverse image of the
closed set on its right. Lemma~\ref{lem:border} therefore implies
$\Phi(\C^{45})\subseteq X_7$. All high-degree coefficients are retained
in this argument and vanish by continuity at the limit.

\paragraph{A nonsingular integer Jacobian.}
Set every parameter to zero except
\begin{equation}\label{eq:point}
 \alpha=b_2=d_3=d_5=d_6=z_8=1.
\end{equation}
Then
\[
 Q=x^4+x^3,\quad R=x^5,\quad P=x^{12},\quad F=x^{17},
 \quad G=x^{22}+x^{19},\quad H=G(G+Q+P+F).
\]
For an explicit derivative table at this point, set
$B=R+x^2$, $K=2G+Q+P+F$, $L=KB+G$. The Jacobian columns, viewed as
polynomials, in the order \eqref{eq:parameters}, are
\begin{center}
\begin{tabular}{@{}ll@{}}
\toprule
Parameter group & Derivative polynomials\\
\midrule
$\alpha$ & $Gx^3$\\
$\beta$ & $(LP+KF)x^3$\\
$\xi_j$ ($j\in J$) & $(LR+G)x^j$\\
$\mathbf u$ & $LR\,L_4$\\
$\mathbf v$ & $LP\,L_3$\\
$\mathbf a$ & $KB\,L_5$\\
$\mathbf b$ & $KF\,L_3$\\
$\mathbf c$ & $(G+Q+P+F)L_6$\\
$\mathbf d$ & $G L_6$\\
$\mathbf z$ & $(1,x,x^2,Q,R,P,F,G,H)$\\
\bottomrule
\end{tabular}
\end{center}
With coefficient rows $0,\ldots,44$, exact elimination gives
\begin{equation}\label{eq:det}
 \det D\Phi=256\ne0.
\end{equation}
The complete integer matrix is supplied with the independent verifier.
It is reconstructed by simultaneous first-order polynomial arithmetic and
compared entrywise with this derivative table. All coefficients are retained.
The entries lie between $0$ and $6$; modular determinants are $1,54,256$
modulo $3,101,1009$, respectively. A separate whole-circuit interpolation
check recovered from the original conversation is also supplied.

The complex inverse function theorem now shows that the image of $\Phi$
is Zariski dense in $\V{44}$. Since $X_7$ is closed and contains that
image, it contains $\V{44}$. This proves Theorem~\ref{thm:main}.

\section{An additional upper bound}
\begin{corollary}\label{cor:bounds}
Using Theorem 10 of Jarlebring and Lorentzon \cite{JL}, one has
$44\le d_7\le47$.
\end{corollary}
\begin{proof}
All seven-product computations can be parameterized by one affine space:
at step $k$, multiply two arbitrary linear combinations of $1,x$ and the
$k-1$ previous product outputs, and finally take an arbitrary linear
combination of $1,x$ and the seven product outputs. Unused products can be
padded by zero products. Thus $X_7$ is the closure of the image of an
irreducible affine space, and is irreducible. The cited dimension theorem
gives $\dim X_7=49$. If $\V{48}\subseteq X_7$, this closed irreducible
subspace has the same dimension as $X_7$, so it must equal $X_7$.
But seven successive squarings compute $x^{128}$, contradicting that
equality. Therefore $d_7\le47$, and the lower bound follows from
Theorem~\ref{thm:main}.
\end{proof}
Determining which of $44,45,46,47$ equals $d_7$ remains open in this
contribution. That further question is not needed to refute the retained
MF-14 equality. The lower bound and negative resolution do not depend on
the external dimension theorem.

\section{Reproduction and evidence scope}
The accompanying submission README records exact commands, source hashes,
machine-readable certificates and two separate review reports. The symbolic
border verifier and the coefficient-map verifier use exact integer/rational
arithmetic. Their checks validate the finite identities above; the analytic
and algebraic closure arguments have been informally reviewed by separate
agents. No Lean proof or kernel-checking claim is made.

\appendix
\section{Explicit high-degree coefficient identities}\label{app:cofactors}
Let $c_i$ be the signed cofactor multiplying $f_i$ in the expansion of
\eqref{eq:E} along its last row. The first, second and fifth polynomials
$f_i$ have degree at most ten, so only $c_3QR_s+c_4R_s^2$ contributes
above degree ten. Direct expansion of the $4\times4$ cofactors gives
$c_3=s^2 C(s)$ and $c_4=s^3 D(s)$, where
\begin{align*}
 C(s)={}&-1-3\alpha\gamma s
 +(-2\alpha^3+9\alpha^2\gamma^2-6\alpha\eta-2\eta\gamma^2)s^2\\
 &+(-4\alpha^4\gamma+3\alpha^3\gamma^3+12\alpha^2\eta\gamma)s^3\\
 &+(-8\alpha^6-4\alpha^5\gamma^2-10\alpha^4\eta)s^4,\\
 D(s)={}&2\gamma+(3\alpha^2-\alpha\gamma^2-\eta)s
             -8\alpha^3\gamma s^2+9\alpha^5s^3.
\end{align*}
Consequently the exact high-degree coefficients are
\begin{align*}
 [x^{11}]E_s&=s^4\{3\alpha C+D[2\gamma+
        (6\alpha^2+2\eta+2\alpha\gamma^2)s+2\alpha^3\gamma s^2]\},\\
 [x^{12}]E_s&=s^4\{C+sD[(6\alpha+\gamma^2)+6\alpha^2\gamma s+\alpha^4s^2]\},\\
 [x^{13}]E_s&=s^5D(2+6\alpha\gamma s+4\alpha^3s^2),\\
 [x^{14}]E_s&=s^6D(2\gamma+6\alpha^2s),\\
 [x^{15}]E_s&=4\alpha s^7D,\qquad [x^{16}]E_s=s^7D.
\end{align*}
These formulas directly give \eqref{eq:lead11}--\eqref{eq:lead12} and all
required divisibilities, including special choices where leading coefficients
vanish. They are checked independently against the full determinant expansion
by the border verifier.

\begin{thebibliography}{9}
\bibitem{JL} E. Jarlebring and G. Lorentzon,
\emph{The polynomial set associated with a fixed number of matrix-matrix
multiplications}, arXiv:2504.01500v3, Theorem 10 and Conjecture 13.
\url{https://arxiv.org/html/2504.01500v3}.
\bibitem{Colbrook} M. J. Colbrook, MF-14 submission, 11 September 2026,
Theorem 1 and equations (1)--(2). Archived in
\href{https://github.com/ajt60gaibb/OpenProblemsInNLA/blob/main/references/colbrook-matrix-functions-2026-09-11/manuscripts/MF-14.tex}{the OpenProblemsInNLA manuscript archive}.
\bibitem{chat} Author-supplied ChatGPT conversation, accessed 17 September
2026. Degree-$44$ construction and subsequent independent-method validation
within the same conversation.
\url{https://chatgpt.com/share/6aaba6d9-3424-83eb-b278-21e3b917a83b}.
\end{thebibliography}
\end{document}
